\section{Introduction}
\label{sec:introduction}

For most of the history of software engineering, source code has been the scarce, durable, and expensive artefact of the discipline. The cost of producing a working implementation dominated every other concern, and the surrounding apparatus---requirements, designs, tests, and documentation---existed largely to amortize and protect that investment. The recent maturation of large language models (LLMs) for code generation has begun to invert this economic order. When a sufficiently capable agent can regenerate an implementation from a description in seconds and at negligible marginal cost, the implementation ceases to be a treasured asset and becomes closer to a disposable build output, akin to an object file recompiled on demand. What then becomes scarce and durable is not the code but the \emph{specification of intent}: the precise statement of what the code must do, against which any number of cheap regenerations can be measured. This article takes the economic inversion seriously and asks what a software lifecycle reorganised around a machine-checkable specification as its source of truth would have to look like.

\subsection{The rise---and the defects---of spec-driven, agentic lifecycles}
\label{subsec:intro-agentic}

The industry has not failed to notice the inversion. A cluster of emerging methodologies and tools now place a ``specification'' at the centre of an AI-assisted lifecycle in which agents generate, test, and revise code: GitHub's Spec Kit~\cite{githubspeckit}, Amazon's AI-Driven Development Life Cycle (AI-DLC)~\cite{awsaidlc}, and the Kiro spec-driven integrated development environment~\cite{awskiro} are representative, and position papers have begun to argue that LLMs are reshaping the software development life cycle wholesale~\cite{ozkaya2023llmse}. These approaches are genuine advances over unstructured prompting because they recognise that intent, not the keystroke-level act of typing code, is the durable artefact. Yet they share three structural defects that this article seeks to remedy.

First, the specification is expressed in \emph{natural language} (NL). An NL specification cannot be mechanically compiled into code, only interpreted, so the central ``compilation'' step from intent to implementation is fundamentally unverifiable; there is no formal object against which the result can be checked. Second, the loop is \emph{open}. The agent generates an artefact and a human, or at best another model, eyeballs it; there is no machine gate that an output must pass before it is admitted, and so regressions and subtle defects propagate silently. Third, where automated verification is present at all, it is typically \emph{circular} or self-grading: the same model that writes the code also writes its tests or serves as its oracle, so the acceptance signal is drawn from the very distribution that produced the candidate. The reliability literature has repeatedly shown that such signals are fragile---benchmark suites are often too weak to detect plausible-but-wrong code, and tightening them lowers reported success rates sharply~\cite{liu2023evalplus}---and that LLM code generation is prone to confident hallucination that survives superficial review~\cite{zhang2025hallucination}. The hazard is acute precisely when an LLM is asked to generate the oracle it will be judged against~\cite{liu2026oracle}.

\subsection{The proposed pipeline: a verification spine for AI-native development}
\label{subsec:intro-proposal}

The thesis of this article is that the missing element in these lifecycles is a \emph{machine-checkable formal-specification contract layer} together with a \emph{verification loop} that closes the gap between intent and implementation. The contract layer supplies the verification spine that agentic methodologies currently lack. We do not propose a new model or a new prompting trick; we propose a reference architecture that re-couples two intellectual lineages that the LLM era has split apart---classical specification-driven program synthesis, with its closed, counterexample-driven loop~\cite{manna1980deductive,solarlezama2008sketching}, and description-driven generation, with its cheap NL intent~\cite{chen2021codex,jimenez2024swebench}---and we validate the architecture by instantiating it.

The pivot of the architecture is the formal contract, which mediates between the ambiguous human world above it and the precise machine world below it. Above the pivot lies \emph{elaboration}: the transformation of NL intent into a formal contract, confirmed by a human. Below the pivot lie two operations on a now-precise artefact: \emph{synthesis}, in which an agent produces code from the contract, and \emph{verification}, in which a deductive verifier checks the code against the contract over \emph{all} inputs rather than a finite sample. Two principles distinguish this organisation from the spec-driven tools surveyed above. The first is the \emph{dual role of the specification}: the same formal contract serves simultaneously as the \emph{input} the agent builds against and the \emph{oracle} the agent's output is verified against. Because the specification plays both roles, the cost of authoring it---historically the bottleneck that crippled formal methods---is the single cost that must be eliminated, and automated inference is what eliminates it. The second is the \emph{independent-authority constraint}: the contract that code is verified against must not be authored by the same agent instance or role that writes the code, on pain of reintroducing the self-grading circularity that the architecture exists to remove. The elaborator role is maximally faithful and strict and never sees the implementation; the implementer role merely satisfies the contract.

Establishing a human-owned oracle is harder than it appears, and naming the difficulty precisely is one of this article's contributions. We term it the \emph{ratification gap}: the specification is the oracle and must therefore be correct and human-owned, yet the human reached for NL in the first place precisely because they cannot or will not write formal contracts, so asking them to confirm formal text they cannot read makes their authority a fiction. The architecture resolves this by a discipline of \emph{ratifying behaviour, not logic}---rendering a candidate contract as concrete, checkable examples (``given $x=-1$, this contract says the result is $0$; is that correct?'') whose ratified instances become \emph{pinned oracles} the contract must satisfy---and by \emph{Socratic, decision-surfacing elaboration}, in which the agent surfaces only the boundary decisions it was forced to resolve (null and empty inputs, overflow, ordering, duplicates, concurrency, error-versus-exception) rather than the full formal text. To guard against an elaborator that quietly weakens the contract to make it trivially satisfiable---``sort'' silently relaxed to ``is a permutation,'' which the identity function satisfies---a second, independent \emph{adversarial elaboration} agent searches for scenarios consistent with the NL intent but violating the proposed contract, the elaboration-layer dual of counterexample-guided synthesis, supplemented by vacuity and non-triviality checks that repurpose our own prior ``ensures true'' punt detection as a weakening signal. Throughout, every obligation in a contract is tagged with its \emph{provenance} (human-ratified, inherited-default, inferred-unconfirmed, example-pinned, ticket-stated, doc-stated, or code-inferred) and its \emph{assurance} level (proved, bounded-checked, tested, or unchecked), which is what makes the independent-authority constraint auditable and the word ``verified'' honest. Where intent is not formally specifiable or verification does not terminate, the architecture degrades gracefully from full proof to bounded check to test, recording the level rather than silently dropping the obligation.

\subsection{Brownfield reality and compositional propagation}
\label{subsec:intro-brownfield}

A reference architecture that assumes a clean slate would be of little use, because most software already exists. For pre-existing (brownfield) codebases the contract is \emph{recovered}, not authored, and the existing code occupies a peculiar dual position: it is simultaneously the best available evidence of intent and the prime suspect, because it encodes whatever bugs it contains. We therefore introduce a \emph{characterisation phase} that runs before any user input, triangulating three streams of evidence---the code (what the system does), the issue tickets (why, including known bugs and edge cases), and the documentation (the declared intent)---and reconciling them into a candidate contract with per-clause provenance, flagging the disagreements between streams. It must be stated precisely what such verification proves: checking code against a \emph{code-derived} specification is characterisation, a faithful-description baseline and a toolchain sanity check; it is nearly circular and does \emph{not} prove the absence of bugs. Bugs surface instead as \emph{disagreements}---the code violating its own inferred pattern, or the code contradicting ticket- or doc-derived intent---so the characterisation phase cannot be gated on ``zero errors''; finding code-versus-intent violations is the success outcome, and its output is a characterisation report listing the recovered specification, the candidate bugs, and what could not be verified, which becomes the opening agenda of the human dialogue. This motivates a \emph{reordering} of the lifecycle loop: a first, retroactive pass establishes a trusted baseline on existing code before any change is made, and only subsequent, prospective passes implement new features. Because whole legacy applications cannot be formalised at once, the architecture adopts a \emph{frontier} strategy, specifying only the module under change, the public API surface, and the modules referenced by active tickets, and treating the remainder as unspecified under weak inherited defaults.

It is at these module boundaries that the thesis programme's central mechanism plugs in. Because formal contracts compose, a change to a callee's specification \emph{recomputes} the proof obligations of its callers and flags which callers break, \emph{without running them}. Internal refactoring and dependency or library version upgrades become the same operation, and the perennially expensive question ``which of my call sites does this upgrade break behaviourally?'' becomes a first-class, machine-answered lifecycle event---a blast-radius or behavioural-compatibility analysis. This lifts the compositional refinement established in our prior work~\cite{edmonds_article3} from a property of the inference engine into a lifecycle capability, and it directly serves the broader programme goal of supporting formal-methods adoption and library version compatibility.

\subsection{Relation to prior work in this programme}
\label{subsec:intro-prior}

This is the fifth article in a research programme on automated specification inference and its uses. The instrument throughout is the JML-Inferrer, a Java~21 tool built on JavaParser that infers Java Modeling Language (JML)~\cite{leavens2006jml} pre- and postconditions, loop and class invariants, and assignable clauses by abstract-syntax-tree heuristics, with verification by OpenJML extended static checking~\cite{cok2014openjml} discharged through the Z3 SMT solver~\cite{demoura2008z3}, augmented by a custom OpenJML fork that adds recursive function definitions for the \texttt{\textbackslash sum}, \texttt{\textbackslash product}, and \texttt{\textbackslash num\_of} quantifiers. The four prior articles supply the trusted, already-validated components on which the present architecture's proof-of-concept is built: Article~1 of this programme showed that automatically inferred JML specifications materially improve LLM-based unit-test generation~\cite{edmonds_article1}; Article~2 showed that JML specifications can be embedded into compiled Java artefacts and OpenAPI documents at full roundtrip fidelity~\cite{edmonds_article2}; Article~3 showed that heuristic inference admits compositional refinement across method boundaries, a second weakest-precondition pass adding on the order of $18{,}854$ additional well-formed preconditions on the study corpus~\cite{edmonds_article3}; and Article~4 integrated the inference pipeline into continuous-integration and continuous-delivery (CI/CD) systems with diff-only analysis, content-hash caching, and fail-open semantics~\cite{edmonds_article4}. The present article reuses these as trusted building blocks and contributes the lifecycle architecture that binds them together.

\subsection{Research questions and contributions}
\label{subsec:intro-rqs}

The proposed architecture (contribution~(i) below) frames a single overarching
design question---what lifecycle architecture supplies the missing verification
spine for agentic, spec-driven development, and what roles, artefacts, and
invariants must it enforce so that the formal contract can act simultaneously as
input intent and verification oracle without self-grading. That design question
is decomposed into six research questions, each operationalised in the evaluation
design (\cref{sec:evaluation_plan}).

\begin{itemize}
  \item \textbf{RQ1 (Ratification effort).} Can ratifying behaviour through pinned-oracle examples let a non-expert legitimately own a formal contract, and at what human cost relative to confirming formal text?
  \item \textbf{RQ2 (Decision surfacing and adversarial soundness).} Does Socratic boundary-decision surfacing capture the decisions where defects concentrate, and does adversarial elaboration detect specification weakening?
  \item \textbf{RQ3 (Provenance, assurance, and independent authority).} Do per-obligation provenance and assurance tags, together with the independent-authority constraint, make a ``verified'' claim auditable?
  \item \textbf{RQ4 (Synthesis--verification loop).} Does the counterexample-guided loop discharge obligations on the method classes the inferrer supports, and how does tiered assurance degrade where it does not?
  \item \textbf{RQ5 (Brownfield characterisation).} For pre-existing codebases, can multi-source triangulation of code, tickets, and documentation recover a provenance-tagged contract and surface genuine code-versus-intent discrepancies as candidate bugs on a real three-stream corpus?
  \item \textbf{RQ6 (Compositional propagation).} Can compositional contracts make behavioural version-compatibility a first-class lifecycle event, identifying which callers a callee-specification change breaks without executing them?
\end{itemize}

\noindent The corresponding contributions are: (i) a layered reference architecture with the formal contract as the pivot between the human and machine worlds, the dual-role specification, and the independent-authority constraint (\cref{sec:architecture} develops it); (ii) an elaboration discipline that resolves the ratification gap by ratifying behaviour rather than logic, with adversarial elaboration and provenance- and assurance-tagged obligations; (iii) a brownfield characterisation phase and loop reordering that treats disagreement detection, not error-freedom, as the success criterion; (iv) the promotion of compositional refinement to a behavioural-compatibility, blast-radius capability; and (v) a proof-of-concept instantiation that reuses the programme's validated assets as trusted components---the JML-Inferrer for draft contracts, the OpenJML fork and Z3 as the oracle that never grades itself---instantiated as a partially built proof of concept on Apache Commons Lang, in which the synthesis--verification loop is exercised against fixed contracts (milestone M1) while contract generation and conversational ratification (milestones M2--M3) are specified as a staged build plan. The corpus uniquely couples code with a public issue tracker (the ``LANG-\#\#\#\#'' Apache JIRA project) and rich Javadoc, yielding a genuine three-stream dataset on which real code-versus-intent discrepancies can be exhibited.

We are deliberate about what is and is not claimed. This article \emph{proposes and partially instantiates} a reference architecture; it argues for the architecture's feasibility and realism rather than a controlled effect size. Several limitations are load-bearing and stated openly: inferred specifications skew weak because they describe what code does rather than what it should, making ratification real labor that we mitigate with inherited default contracts and API-surface prioritisation; inferred specifications are sometimes over-strong and yield false candidate bugs, which is why every discrepancy is surfaced as a \emph{candidate} for human adjudication; NL-to-formal translation without a human is imperfect, so the resulting clauses remain unconfirmed and drive the dialogue rather than gating it; linking tickets and documentation to code is noisy and credible only on a curated module; verification scalability and decidability limits, including Z3 timeouts on multiplicative reasoning, array theory, and rich preconditions, are real and known; and, most importantly, a wrong specification verified against is worse than none, which is exactly why the independent-authority constraint and human ratification are not optional decorations but the architecture's structural load-bearers. Full empirical validation on external production codebases and an industrial action-research case study, with the attendant human-research ethics approval, is future work.

\subsection{Roadmap}
\label{subsec:intro-roadmap}

The remainder of the article is structured as follows. \cref{sec:related_work} synthesizes the related work across program synthesis, specification inference, verified code generation, the test-oracle problem, and behavioural version compatibility, positioning the proposed pipeline against each. The subsequent sections present the layered reference architecture and its invariants; the elaboration layer and its resolution of the ratification gap; the brownfield characterisation phase and loop reordering; the synthesis--verification loop and its compositional change-propagation capability; and the proof-of-concept instantiation, including the verify-as-a-service, fixed-contract synthesis, contract-generation, and conversational-ratification milestones evaluated on Apache Commons Lang. The article closes with a discussion of threats to validity, limitations, and the future industrial-validation programme.
